% SGA 6 English synchronization fragment.
% Deliverable range only: current-rescribe idx625--635.
% Printed volume pages 612--622; declared source-PDF pages 615--625.
% French control: sga6_fr_workpass.tex (Claude workpass, source-checked through idx646).
% Scan authority: local sga6.pdf; high-resolution supplement used only to
% disambiguate symbols and source-level typographical defects.

\providecommand{\Isom}{\operatorname{Isom}}
\providecommand{\Coker}{\operatorname{Coker}}
\providecommand{\Ker}{\operatorname{Ker}}
\providecommand{\Supp}{\operatorname{Supp}}
\providecommand{\Ass}{\operatorname{Ass}}

% Current rescribe index 625; printed volume page 612; source-PDF page 615.
There is a canonical monomorphism from \(\mathrm{Deff}\) into the functor
\(\mathrm{Rec}\) of gluing data on \(L\). By definition,
\(\mathrm{Rec}\) is the functor
\(\Isom(p_1^*L,p_2^*L)\), where the
\(p_i:X\times_YX\to X\) are the projections; by \hyperlink{sga6.e12.lem.2-4}{2.4} this functor is
represented by an affine scheme \(I\) of finite presentation over \(S\). It
therefore remains to represent the monomorphism
\[
  i:\mathrm{Deff}\hookrightarrow \Isom(p_1^*L,p_2^*L)
\]
by a morphism of finite presentation, which will necessarily be quasi-affine
(by EGA IV 8.11.2).

Let us note in passing that, if desired, one may restrict
\(\Isom(p_1^*L,p_2^*L)\) by imposing the cocycle condition, but this is not
needed below.

After the base change \(I\to S\), we are thus reduced to the following
problem: if \(L\) is an invertible sheaf on \(X\), endowed with a gluing datum
\[
  \varphi:p_1^*L\xrightarrow{\sim}p_2^*L,
\]
represent the subfunctor of \(S\) on which this datum becomes effective. We
have the following lemma.

\sgaStableTarget{sga6.e12.lem.4-3}\textbf{Lemma 4.3.} Under the hypotheses of \hyperlink{sga6.e12.prop.4-1}{4.1}, let \(L\) be an invertible
sheaf on \(X\), endowed with a gluing datum
\[
  \varphi:p_1^*L\xrightarrow{\sim}p_2^*L.
\]
This datum is effective if and only if the following five conditions hold.
\begin{enumerate}[label=\arabic*)]
\item For every \(i\geq0\), \(R^if_*(L)\) is flat over \(S\).
\item For every \(j\geq0\), \(R^jg_*(p_1^*L)\) is flat over \(S\).
\item The cokernel of
\[
  \pi_1-\pi_2:f_*(L)\longrightarrow
  g_*(p_1^*L)\simeq g_*(p_2^*L)
\]
is flat over \(S\), where the latter isomorphism is \(g_*(\varphi)\).
\item The kernel \(N\) of
\(\pi_1-\pi_2:f_*(L)\to g_*(p_1^*L)\) is invertible on \(Y\).
\item The composite morphism
\[
  f^*N\longrightarrow f^*f_*L\longrightarrow L
\]
is an isomorphism.
\end{enumerate}

\newpage
% Current rescribe index 626; printed volume page 613; source-PDF page 616.
Indeed, these conditions are sufficient, since (4) and (5) alone imply
effectivity. Conversely, let \(M\) be invertible on \(Y\), and suppose that
there is an isomorphism \(f^*M\simeq L\) compatible with the gluing datum. We
may cover \(Y\) by open subsets \(V\) over which \(M\) is isomorphic to
\(\mathcal O_Y\). The descent datum on \(L\), restricted over
\(f^{-1}(V)\), is therefore isomorphic to the trivial descent datum on
\(\mathcal O_X|_{f^{-1}(V)}\). Finally, conditions (1)--(5) are local on
\(Y\), and the hypotheses of \hyperlink{sga6.e12.prop.4-1}{4.1} show that they hold for
\(L=\mathcal O_X\); hence they hold for \(L\).

It remains to represent conditions (1)--(5), one after another, by
monomorphisms of finite presentation. For (1) and (2), we apply the following
lemma.

\sgaStableTarget{sga6.e12.lem.4-4}\textbf{Lemma 4.4.} Let \(S\) be a scheme, let
\(\lambda:X\to S\) and \(\mu:Y\to S\) be proper morphisms of finite
presentation, and let \(f:X\to Y\) be an \(S\)-morphism. Let \(F\) be a
sheaf on \(X\), of finite presentation and flat over \(S\). For an integer
\(n\), let \(S_n\) be the subfunctor of \(S\) on which \(R^if_*(F)\) is
flat over \(S\) for every \(i\geq n\); more precisely, for a morphism
\(\phi:T\to S\),
\[
  S_n(T)=
  \begin{cases}
    \{\phi\},&\text{if \(R^i(f_T)_*(F_T)\) is flat for every \(i\geq n\),}\\
    \varnothing,&\text{otherwise.}
  \end{cases}
\]
Then \(S_n\hookrightarrow S\) is representable by a surjective monomorphism
of finite presentation; moreover, for \(i\geq n\), the formation of
\(R^i(f_{S_n})_*(F_{S_n})\) commutes with base changes \(T\to S_n\). In
addition, \(S_n\) coincides with the functor \(S'_n\) on which the functors
\[
  M\longmapsto
  R^i(f_V)_*\bigl(F_V\otimes\lambda_V^*M\bigr),
  \qquad i\geq n,
\]
are exact, where \(M\) ranges over the quasi-coherent
\(\mathcal O_V\)-modules, for every affine open subset \(V\) of \(S\).

Indeed, the assertion is local on \(S\), so we may suppose that \(S\) is
quasi-compact. The relative dimension of \(X/Y\) is then bounded, say by
\(r\), and, for \(n>r\),
\[
  R^nf_*(F_V\otimes\lambda_V^*M)=0.
\]
The assertion therefore holds for \(n>r\).

\newpage
% Current rescribe index 627; printed volume page 614; source-PDF page 617.
Proceeding by descending induction, suppose it is known for \(n+1\). After
replacing \(S\) by \(S_{n+1}=S'_{n+1}\), we may suppose that
\(R^if_*(F)\) is flat over \(S\) for \(i\geq n+1\). Since \(F\) is flat
over \(S\), EGA III 6.9.8 then shows that the formation of \(R^nf_*(F)\)
commutes with every base change. Consequently, the representability of
\(S_n\) follows from the next lemma (see J. P. Murre,
\emph{Representation of unramified functors}, S\'eminaire Bourbaki, May
1965, p.~294-11, Theorem~2).

\sgaStableTarget{sga6.e12.lem.4-5}\textbf{Lemma 4.5.} Let \(X\to S\) be a proper morphism of finite
presentation, and let \(F\) be a sheaf of finite presentation on \(X\). Then
the subfunctor \(S_F\) of \(S\) on which \(F\) is flat over \(S\) is
representable by a surjective monomorphism of finite presentation (which, of
course, is not an immersion in general).

To represent condition (3) of \hyperlink{sga6.e12.lem.4-3}{4.3}, note that, by (1) and (2), the formation
of \(f_*(L)\) and \(g_*(p_1^*L)\) commutes with the base changes
\(I\to S\); hence so does the formation of
\(\Coker(\pi_1-\pi_2)\). We then apply 4.5.

To represent (4), note that, by (3), and because \(f_*(L)\) and
\(g_*(p_1^*L)\) are flat and commute with base change,
\[
  N=\Ker(\pi_1-\pi_2)
\]
also commutes with base change. We then apply the following lemma.

\sgaStableTarget{sga6.e12.lem.4-6}\textbf{Lemma 4.6.} Let \(\lambda:X\to S\) be a proper morphism of finite
presentation, and let \(u:F\to G\) be a homomorphism of
\(\mathcal O_X\)-modules of finite presentation, with \(G\) flat over
\(S\). Then the subfunctor \(S_1\) of \(S\) on which \(u\) is an
isomorphism is representable by an open subscheme of \(S\), of finite
presentation.

As usual, we reduce to the case where \(S\) is noetherian, using EGA IV 8
and 11.2.6; this relieves us of the need to consider finite presentation.
First let us represent the functor on which \(u\) is an epimorphism,

\newpage
% Current rescribe index 628; printed volume page 615; source-PDF page 618.
that is, on which \(\Coker u\) is zero. Since the formation of
\(\Coker u\) commutes with base change, this functor is represented by the
open subset
\[
  S-\lambda\bigl(\Supp(\Coker u)\bigr)
\]
of \(S\). We may therefore suppose that \(u\) is already an epimorphism.
Using now the hypothesis that \(G\) is flat over \(S\), we find that the
formation of \(\Ker u\) commutes with base change; consequently, the functor
under consideration is represented by the open subset
\[
  S-\lambda\bigl(\Supp(\Ker u)\bigr)
\]
of \(S\).

\newpage
% Current rescribe index 629; printed volume page 616; source-PDF page 619.
\sgaStableTarget{sga6.e13.expose}\section*{Expos\'e XIII. Finiteness Theorems for the Picard Functor}
\begin{center}
by S. Kleiman
\end{center}

\sgaStableTarget{sga6.e13.sec.1}\subsection*{1. The \texorpdfstring{\((b)\)}{(b)}-sheaves}

Let \(k\) be an algebraically closed field and \(X\) a projective
\(k\)-scheme, endowed with an ample invertible sheaf \(\mathcal O_X(1)\).

\sgaStableTarget{sga6.e13.def.1-1}\textbf{Definition 1.1 (Mumford).} Let \(m\in\mathbb Z\). A coherent sheaf
\(F\) on \(X\) is said to be \emph{\(m\)-regular} if, at every point of the
support of \(F\), \(\mathcal O_X(1)\) is generated by its global sections
and if, for every \(q\geq1\),
\[
  H^q(F(m-q))=0.
\]

\sgaStableTarget{sga6.e13.lem.1-2}\textbf{Lemma 1.2.} Let
\[
  0\longrightarrow F(-1)\longrightarrow F\longrightarrow G\longrightarrow0
\]
be an exact sequence of coherent sheaves on \(X\).\footnote{Here, as below
in analogous situations, it is understood that \(F(-1)\to F\) is defined
by tensoring with a given section of \(\mathcal O_X(1)\).} If \(F\) is
\(m\)-regular, then so is \(G\).

Indeed, we have the corresponding exact sequence
\[
  H^q(F(m-q))\longrightarrow H^q(G(m-q))
  \longrightarrow H^{q+1}(F(m-q-1)).
\]

\sgaStableTarget{sga6.e13.prop.1-3}\textbf{Proposition 1.3.} If a sheaf \(F\) on \(X\) is \(m\)-regular,
then, for every \(n\geq m\),
\begin{enumerate}[label=(\roman*)]
\item \(F\) is \(n\)-regular.
\item The map
\[
  H^0(F(n))\otimes H^0(\mathcal O_X(1))
  \longrightarrow H^0(F(n+1))
\]
is surjective.
\item \(F(n)\) is generated by \(H^0(F(n))\).
\end{enumerate}

\newpage
% Current rescribe index 630; printed volume page 617; source-PDF page 620.
The proof is by induction on the dimension \(s\) of the support of \(F\); if
\(s=-1\), the assertion is trivial. Choose a section
\(\sigma\in H^0(\mathcal O_X(1))\) which vanishes at no point of the finite
set \(\Ass(F)\). Then \(\sigma\) defines an exact sequence
\[
  0\longrightarrow F(-1)\longrightarrow F\longrightarrow G\longrightarrow0,
\]
and \(G\) is \(m\)-regular by \hyperlink{sga6.e13.lem.1-2}{1.2}, with support of dimension \(s-1\).

Consider the corresponding exact sequence
\[
  H^q(F(n-q-1))\longrightarrow H^q(F(n-q))
  \longrightarrow H^q(G(n-q)).
\]
By the induction hypothesis, \(H^q(G(n-q))=0\) for \(n\geq m\), so the
vanishing of \(H^q(F(n-q-1))\) implies that of \(H^q(F(n-q))\). But
\(H^q(F(m-q))=0\), which proves (i).

To prove (ii), we chase through the following commutative diagram:
\[
\begin{tikzcd}[column sep=1.7em,row sep=2.4em]
  & & H^0(F(n))\otimes H^0(\mathcal O_X(1))
      \arrow[rr] \arrow[d]
  & & H^0(G(n))\otimes H^0(\mathcal O_X(1))
      \arrow[r] \arrow[d]
  & 0 \\
  0 \arrow[r]
  & H^0(F(n)) \arrow[ur] \arrow[r]
  & H^0(F(n+1)) \arrow[rr]
  & & H^0(G(n+1)) \arrow[d] \\
  & & & & 0
\end{tikzcd}
\]
Here the oblique arrow is defined by
\(\varphi\mapsto\varphi\otimes\sigma\); the horizontal rows are exact, the
first by the \(n\)-regularity of \(F\) and by (i), already proved, and the
second column is exact by the induction hypothesis.

Finally, (iii) follows from the commutative diagram
\[
\begin{tikzcd}[column sep=large,row sep=large]
H^0(F(n))_X\otimes H^0(\mathcal O_X(1))_X
  \arrow[r,"\alpha_n"]
  \arrow[d]
& H^0(F(n)(1))_X
  \arrow[d,"\beta_{n+1}"] \\
{[H^0(F(n))_X](1)}
  \arrow[r,"\beta_n\otimes\mathrm{id}"]
& F(n)(1).
\end{tikzcd}
\]

\newpage
% Current rescribe index 631; printed volume page 618; source-PDF page 621.
Here, if \(M\) is a vector space over \(k\), we write \(M_X\) for
\(M\otimes_k\mathcal O_X\). For \(n\geq m\), \(\alpha_n\) is surjective
by (ii); therefore, if \(\beta_{n+1}\) is surjective, so is \(\beta_n\). By
Serre's theorem, \(\beta_{n+1}\) is surjective for \(n\gg0\); this proves
(iii).

\sgaStableTarget{sga6.e13.prop.1-4}\textbf{Proposition 1.4.} Let
\[
  0\longrightarrow F(-1)\longrightarrow F\longrightarrow G\longrightarrow0
\]
be an exact sequence on \(X\). If \(G\) is \(m\)-regular, then:
\begin{enumerate}[label=(\roman*)]
\item \(H^q(F(n))=0\) for \(q\geq2\) and \(n\geq m-q\).
\item \(h^1(F(n-1))\geq h^1(F(n))\) for \(n\geq m-1\).
\item \(H^1(F(n))=0\) for
\(n\geq(m-1)+h^1(F(m-1))\).
\end{enumerate}
In particular, \(F\) is \([m+h^1(F(m-1))]\)-regular.

Indeed, in the exact sequence
\[
  H^{q-1}(G(n))\longrightarrow H^q(F(n-1))
  \longrightarrow H^q(F(n))\longrightarrow H^q(G(n)),
\]
the two outer groups are zero for \(q\geq2\) and
\(n\geq m-(q-1)\), by \hyperlink{sga6.e13.prop.1-3}{1.3(i)}. Thus
\[
  H^q(F(m-q))\xrightarrow{\sim}H^q(F(m-q+1))
  \xrightarrow{\sim}H^q(F(m-q+2))\xrightarrow{\sim}\cdots;
\]
but, by Serre's theorem, \(H^q(F(n))=0\) for \(n\gg0\), proving (i).

For \(n\geq m-1\), consider the exact sequence
\[
  0\to H^0(F(n-1))\to H^0(F(n))
  \xrightarrow{\alpha_n}H^0(G(n))\to H^1(F(n-1))
  \to H^1(F(n))\to0.
\]
It implies (ii). Now suppose that \(\alpha_n\) is surjective, and consider
the commutative diagram
\[
\begin{tikzcd}[column sep=large,row sep=large]
H^0(F(n))\otimes H^0(\mathcal O_X(1))
  \arrow[r,"\alpha_n\otimes\mathrm{id}"]
  \arrow[d]
& H^0(G(n))\otimes H^0(\mathcal O_X(1))
  \arrow[r]
  \arrow[d,"\beta_n"]
& 0 \\
H^0(F(n+1))
  \arrow[r,"\alpha_{n+1}"]
& H^0(G(n+1)).
\end{tikzcd}
\]

\newpage
% Current rescribe index 632; printed volume page 619; source-PDF page 622.
If \(n\geq m\), then \(\beta_n\) is surjective by \hyperlink{sga6.e13.prop.1-3}{1.3(ii)}, and hence
\(\alpha_{n+1}\) is surjective. Consequently,
\[
  H^1(F(n-1))\xrightarrow{\sim}H^1(F(n))
  \xrightarrow{\sim}\cdots\xrightarrow{\sim}0.
\]
On the other hand, if \(\alpha_n\) is not surjective, then
\(h^1(F(n-1))>h^1(F(n))\). Thus the function
\(n\longmapsto h^1(F(n))\) decreases \emph{strictly}, starting from
\(h^1(F(m-1))\), until it becomes zero; this proves (iii).

\sgaStableTarget{sga6.e13.def.1-5}\textbf{Definition 1.5.} Let \(F\) be a coherent sheaf on \(X\), let
\(r\geq\dim\Supp(F)\), and let \((b)=(b_0,\ldots,b_r)\) be a sequence of
\(r+1\) integers. Then \(F\) is said to be a \emph{\((b)\)-sheaf} if it
satisfies either of the following two equivalent sets of conditions.

\noindent\sgaStableTarget{sga6.e13.para.1-5-1}\textbf{(1.5.1)}
\begin{enumerate}[label=(\roman*)]
\item At every point of the support of \(F\), \(\mathcal O_X(1)\) is
generated by \(H^0(\mathcal O_X(1))\).
\item \(h^0(F(-1))\leq b_0\).
\item If \(r\geq1\), there is a section
\(\sigma\in H^0(\mathcal O_X(1))\) which defines an exact sequence
\[
  0\longrightarrow F(-1)\longrightarrow F\longrightarrow G\longrightarrow0
\]
such that \(G\) is a \((b_1,\ldots,b_r)\)-sheaf.
\end{enumerate}

\noindent\sgaStableTarget{sga6.e13.para.1-5-2}\textbf{(1.5.2)} There is an \(F\)-regular sequence of \(r\)
sections \(\sigma_1,\ldots,\sigma_r\in H^0(\mathcal O_X(1))\) such that,
for \(i=0,\ldots,r\), if \(F_i\) denotes the restriction of \(F\) to the
zero scheme of \(\sigma_1,\ldots,\sigma_i\), then
\[
  h^0(F_i(-1))\leq b_i.
\]

\sgaStableTarget{sga6.e13.prop.1-6}\textbf{Proposition 1.6.} Let \(F\) be a \((b)\)-sheaf on \(X\). Then:
\begin{enumerate}[label=(\roman*)]
\item For every sufficiently general sequence \(\sigma\) of \(r\) sections
\(\sigma_1,\ldots,\sigma_r\in H^0(\mathcal O_X(1))\), if
\(F_{\sigma,i}\) denotes the restriction of \(F\) to the zero scheme of
\(\sigma_1,\ldots,\sigma_i\), then
\(h^0(F_{\sigma,i}(-1))\leq b_i\), for \(0\leq i\leq r\).
\item Every subsheaf \(G\) of \(F\) is a \((b)\)-sheaf.
\end{enumerate}

\newpage
% Current rescribe index 633; printed volume page 620; source-PDF page 623.
Indeed, let \(S'=\mathbb A_k^N\) be the affine space whose \(k\)-points
correspond to the sequences \(\sigma\), and let \(T\) be the open subset of
\(S'\) corresponding to the \(F\)-regular sequences. For fixed \(i\), the
sheaves \(F_{\sigma,i}(-1)\) corresponding to the \(k\)-points of \(T\) are
contained in a flat family over \(T\), and \(T\) is nonempty by hypothesis.
Thus (i) follows from upper semicontinuity of the function
\[
  \sigma\longmapsto h^0(F_{\sigma,i}(-1))
\]
(EGA III 7.7.5). Finally, since every sufficiently general sequence is
\(G\)-regular and is such that \(G_{\sigma,i}\) is a subsheaf of
\(F_{\sigma,i}\), for every \(i\), (ii) follows from (i).

\sgaStableTarget{sga6.e13.lem.1-7}\textbf{Lemma 1.7.} Let
\[
  0\longrightarrow F(-1)\longrightarrow F\longrightarrow G\longrightarrow0
\]
be an exact sequence on \(X\). If the Hilbert polynomial of \(F\) is written
in the form
\[
  \chi(F(n))=\sum_{i=0}^{r}a_i\binom{n+i}{i},
\]
then
\[
  \chi(G(n))=\sum_{i=0}^{r-1}a_{i+1}\binom{n+i}{i}.
\]
Indeed,
\[
\begin{aligned}
  \chi(G(n))
  &=\chi(F(n))-\chi(F(n-1))\\
  &=\sum_{i=0}^{r}a_i
    \left[\binom{n+i}{i}-\binom{n-1+i}{i}\right]\\
  &=\sum_{i=1}^{r}a_i\binom{n+i-1}{i-1},
\end{aligned}
\]
which proves the assertion.

\sgaStableTarget{sga6.e13.prop.1-8}\textbf{Proposition 1.8.} Let \(F\) be a \((b)\)-sheaf on \(X\), let
\(s=\dim\Supp(F)\), and let
\[
  \chi(F(n))=\sum_{i=0}^{s}a_i\binom{n+i}{i}
\]
be the Hilbert polynomial of \(F\). Then:
\begin{enumerate}[label=(\roman*)]
\item For \(n\geq-1\),
\[
  h^0(F(n))\leq\sum_{i=0}^{s}b_i\binom{n+i}{i}.
\]
\item \(a_s\leq b_s\), and \(F\) is also a
\((b_0,\ldots,b_{s-1},a_s)\)-sheaf.
\end{enumerate}
Indeed, we argue by induction on \(s\). If \(s=0\), the assertion is clear,
since \(a_0=h^0(F)=h^0(F(-1))\leq b_0\). If \(s\geq1\), we have an exact
sequence
\[
  0\longrightarrow F(-1)\longrightarrow F\longrightarrow G\longrightarrow0
\]

\newpage
% Current rescribe index 634; printed volume page 621; source-PDF page 624.
with \(G\) a \((b_1,\ldots,b_r)\)-sheaf and
\(\dim\Supp(G)=s-1\). Then
\[
  h^0(F(n))-h^0(F(n-1))\leq h^0(G(n)),
\]
whereas
\[
  h^0(G(n))\leq
  \sum_{i=0}^{s-1}b_{i+1}\binom{n+i}{i}
\]
by induction on \(s\), and \(h^0(F(-1))\leq b_0\). Assertion (i) follows
by induction on \(n\geq-1\). Finally, by \hyperlink{sga6.e13.lem.1-7}{1.7} and induction on \(s\),
\(a_s\leq b_s\), and \(G\) is a
\((b_1,\ldots,b_{s-1},a_s)\)-sheaf; this proves (ii).

\sgaStableTarget{sga6.e13.def.1-9}\textbf{Definition 1.9.} The following polynomials with nonnegative rational
coefficients, defined recursively, are called the \emph{\((b)\)-polynomials}:
\[
\begin{cases}
P_{-1}=0,\\
P_r(X_0,\ldots,X_r)=P_{r-1}(X_1,\ldots,X_r)
+\displaystyle\sum_{i=0}^{r}X_i
\binom{P_{r-1}(X_1,\ldots,X_r)-1+i}{i}.
\end{cases}
\]

\sgaStableTarget{sga6.e13.rem.1-10}\textbf{Remark 1.10.} For the \((b)\)-polynomials, induction on \(r\geq t\)
immediately gives
\[
  P_r(X_0,\ldots,X_t,0,\ldots,0)=P_t(X_0,\ldots,X_t).
\]

\sgaStableTarget{sga6.e13.thm.1-11}\textbf{Theorem 1.11.} \emph{(The main theorem on \((b)\)-sheaves.)} Let
\(F\) be a \((b)\)-sheaf on \(X\), and let
\[
  \chi(F(n))=\sum_{i=0}^{r}a_i\binom{n+i}{i}
\]
be the Hilbert polynomial of \(F\). Let
\((c)=(c_0,\ldots,c_r)\) be a sequence of integers satisfying
\((c)\geq(b)-(a)\), and let \(m=P_r(c)\) be the integer given by the
\(r\)-th \((b)\)-polynomial \hyperlink{sga6.e13.def.1-9}{(1.9)}. Then \(m\geq0\), and \(F\) is
\(m\)-regular. In particular, if \(s=\dim\Supp(F)\), then \(F\) is
\[
  P_{s-1}(c_0,\ldots,c_{s-1})\text{-regular}.
\]

\newpage
% Current rescribe index 635; printed volume page 622; source-PDF page 625.
Indeed, we argue by induction on \(r\). If \(r=0\), the assertion is clear,
since \(m=0\) and \(F\) is \(n\)-regular for every \(n\). If \(r\geq1\),
there is an exact sequence
\[
  0\longrightarrow F(-1)\longrightarrow F\longrightarrow G\longrightarrow0
\]
in which \(G\) is a \((b_1,\ldots,b_r)\)-sheaf. By \hyperlink{sga6.e13.lem.1-7}{1.7} and the induction
hypothesis, if \(n=P_{r-1}(c_1,\ldots,c_r)\), then \(n\geq0\) and \(G\) is
\(n\)-regular. Hence, by \hyperlink{sga6.e13.prop.1-4}{1.4}, \(F\) is
\([n+h^1(F(n-1))]\)-regular, and \(h^q(F(n-1))=0\) for \(q\geq2\).
Therefore
\[
  h^1(F(n-1))=h^0(F(n-1))-\chi(F(n-1)),
\]
which, by \hyperlink{sga6.e13.prop.1-8}{1.8(i)}, is at most
\[
  \sum_{i=0}^{r}(b_i-a_i)\binom{n-1+i}{i},
\]
since \(b_i\geq0\). Finally, by \hyperlink{sga6.e13.prop.1-3}{1.3(i)}, \(F\) is
\[
  \left[n+\sum_{i=0}^{r}c_i\binom{n-1+i}{i}\right]\text{-regular}.
\]
This proves the first assertion; the second follows from \hyperlink{sga6.e13.prop.1-8}{1.8(ii)} and 1.10.

\sgaStableTarget{sga6.e13.para.1-12}\textbf{1.12.} Let \(S\) be a noetherian scheme, and let \(X\) be an
\(S\)-scheme of finite type. Let \(\mathcal F\) be a family of classes of
coherent sheaves on the fibers of \(X/S\). In other words, for every point
\(s\in S\) and every extension \(K\) of \(k(s)\), coherent sheaves \(F_K\)
are specified on the algebraic scheme \(X_K\); and \(F_K\) and
\(F'_{K'}\) are declared to lie in the same class if there are
\(k(s)\)-homomorphisms from \(K\) and \(K'\) into a common extension
\(K''\) of \(k(s)\) such that
\[
  F_{K''}=F_K\otimes_KK''
  \quad\text{and}\quad
  F'_{K''}=F'_{K'}\otimes_{K'}K''
\]
are isomorphic on \(X_{K''}\). The family \(\mathcal F\) is said to be
\emph{bounded by the coherent sheaf} \(F\) \emph{on}
\(X_T=X\times_ST\), where \(T\) is of finite type over \(S\), if
\(\mathcal F\) is contained in the family of classes of the sheaves
\(F_{k(t)}\), with \(t\in T\). The family \(\mathcal F\) is said to be
\emph{bounded} if such \(T\) and \(F\) exist.

Suppose in addition that \(X/S\) is projective and endowed with a
(relatively) ample invertible sheaf \(\mathcal O_X(1)\). Then
\(\mathcal F\) is called a \((b)\)-\emph{family}, for a sequence of integers
\((b)=(b_0,\ldots,b_r)\), if every class in \(\mathcal F\) is represented
by an \(F_K\), with \(K\) algebraically closed, which is a
\((b)\)-sheaf.
